\section{Hiện thực hệ thống trợ lý AI và Gợi ý hoạt động}

Hệ thống Trợ lý AI của UniConnect được xây dựng dựa trên kiến trúc RAG (Retrieval-Augmented Generation) kết hợp mô hình suy luận ReAct Agent và phân tích ngữ cảnh người dùng theo thời gian thực. Khác với các chatbot hỏi đáp thông thường, Trợ lý AI trong UniConnect có khả năng thấu hiểu lịch trình cá nhân, sở thích của người dùng và dữ liệu bản đồ hoạt động học đường để đưa ra các gợi ý chính xác và có thể thực thi được.

\subsection{Kiến trúc Hybrid Context (Ngữ cảnh hỗn hợp)}

Để sinh câu trả lời và đưa ra gợi ý chất lượng, hệ thống tổng hợp ngữ cảnh từ ba nguồn dữ liệu chính:
\begin{enumerate}
    \item \textbf{Campus Events \& Rules (Tri thức hoạt động):} Thông tin chi tiết về các sự kiện, tiêu chí ngày CTXH, địa điểm tổ chức và quy định công nhận ngày CTXH từ bảng \texttt{activities} thông qua truy vấn tìm kiếm Vector \textit{pgvector}.

    \item \textbf{Personal Schedule \& Interests (Lịch trình và Sở thích):} Dữ liệu thời khóa biểu, khoảng thời gian bận cố định từ bảng \texttt{user\_busy\_slots}, lịch trình tham gia sự kiện đã duyệt trong \texttt{join\_requests} và danh sách từ khóa sở thích (\texttt{interests}) trong hồ sơ người dùng.
    
    \item \textbf{Social Context (Mạng lưới quan hệ):} Các hoạt động có sự tham gia của các Nhóm sinh viên mà người dùng đang là thành viên hoặc những người dùng mà người dùng đang theo dõi (\texttt{user\_follows}), giúp nắm bắt xu hướng phong trào trong cộng đồng.
\end{enumerate}

\subsection{Hiện thực luồng xử lý gợi ý hoạt động cá nhân hóa}

Khi người dùng mở giao diện Bản đồ hoặc yêu cầu Trợ lý AI gợi ý hoạt động, hệ thống thực hiện quy trình xử lý 4 giai đoạn:
\begin{enumerate}
    \item \textbf{Giai đoạn 1 -- Khai thác Vector sở thích:} Hệ thống tính toán vector đại diện cho sở thích của người dùng dựa trên lịch sử hoạt động đã tham gia và danh mục quan tâm trong hồ sơ cá nhân.
    
    \item \textbf{Giai đoạn 2 -- Lọc không gian (Spatial Filtering):} Sử dụng hàm \texttt{ST\_DWithin} của PostGIS để lọc ra các hoạt động đang hoặc sắp diễn ra trong bán kính lân cận khuôn viên trường.

    \item \textbf{Giai đoạn 3 -- Kiểm tra khả dụng thời gian (Calendar Conflict Filtering):} Hệ thống đối soát khung giờ bắt đầu/kết thúc của các hoạt động tiềm năng với bảng \texttt{user\_busy\_slots} và các sự kiện đã đăng ký trong \texttt{join\_requests}. Những sự kiện bị trùng với lịch học cố định hoặc sự kiện đã xác nhận tham gia sẽ được gắn cờ cảnh báo (\texttt{hard\_conflict}) để tránh gợi ý người dùng tham gia vào khung giờ bận.

    \item \textbf{Giai đoạn 4 -- Xếp hạng và Gợi ý (Ranking):} Tổng hợp điểm tương đồng ngữ nghĩa (Cosine Similarity), khoảng cách địa lý và độ ưu tiên từ mạng lưới bạn bè để trả về danh sách Top-K hoạt động tối ưu nhất cho người dùng.
\end{enumerate}

\subsection{Quy trình vận hành RAG Pipeline của Trợ lý AI}

Khi người dùng đặt câu hỏi tự nhiên qua khung chat (ví dụ: \textit{"Thứ Bảy này tôi có thời gian trống vào buổi sáng, có hoạt động tình nguyện nào quanh khu vực trường cấp ngày CTXH không?"}), quy trình xử lý diễn ra như sau:
\begin{enumerate}
    \item \textbf{Text Embedding:} Câu hỏi của người dùng được chuyển đổi thành vector truy vấn 768 chiều thông qua mô hình Gemini Embedding 001 (\texttt{gemini-embedding-001}).
    \item \textbf{Vector Distance Query (pgvector):} Hệ thống thực thi truy vấn tính toán khoảng cách Cosine trực tiếp trên cột \texttt{embedding} (\texttt{Vector(768)}) của bảng \texttt{activities} thông qua toán tử khoảng cách \texttt{cosine\_distance} của thư viện \texttt{pgvector-python}, kết hợp mệnh đề lọc theo trạng thái sự kiện hợp lệ để trích xuất các hoạt động tương đồng nhất.
    \item \textbf{Dynamic Context Injection:} Đóng gói câu hỏi gốc cùng danh sách sự kiện tương đồng trích xuất được và thông tin thời khóa biểu cá nhân của người dùng thành Prompt ngữ cảnh chuẩn hóa.
    \item \textbf{LLM Generation:} Mô hình ngôn ngữ lớn phân tích ngữ cảnh tổng hợp và sinh câu trả lời tư vấn lịch trình tự nhiên, trực quan kèm đường dẫn điều hướng tới trang chi tiết sự kiện trên bản đồ.
\end{enumerate}

\textbf{Ghi chú kỹ thuật về lưu trữ và chỉ mục vector:} Vector đặc trưng ngữ nghĩa được lưu tại cột \texttt{embedding} kiểu dữ liệu \texttt{Vector(768)} của bảng \texttt{activities}. Với quy mô dữ liệu học đường hiện tại, hệ thống thực hiện đối soát khoảng cách Cosine trực tiếp (Exact Nearest Neighbor). Việc cấu hình chỉ mục xấp xỉ đồ thị HNSW (Hierarchical Navigable Small World) trên cơ sở dữ liệu được xác định là giải pháp mở rộng quy mô khi hệ sinh thái sự kiện mở rộng vượt ngưỡng hàng chục nghìn bản ghi.
